\section{The FRB/US model}
\label{sec:model}

\subsection{History and design}

FRB/US was introduced at the Federal Reserve Board in 1996, replacing the MPS model that had been in use since the 1970s \citep{brayton1996guide}. Its design responded to the Lucas-critique era: most behavioural equations are derived from explicit optimisation problems of households, firms and financial-market participants, with expectations appearing explicitly, while retaining the empirical fit and sectoral detail expected of a staff forecasting model. The current public model file specifies 367 variables, of which 284 are endogenous, split between behavioural equations and identities; the remainder are exogenous drivers, policy instruments and switches \citep{fedfrbuspython}.

The model covers the standard expenditure and income sides of the U.S. economy: household consumption of non-durables and services (\texttt{eco}) and durables (\texttt{ecd}); residential (\texttt{eh}) and business fixed investment in equipment (\texttt{epd}, \texttt{epi}) and structures (\texttt{eps}); labour supply, hours (\texttt{lhp}) and unemployment (\texttt{lur}); a detailed price--wage block built around trend unit labour costs, with core PCE inflation (\texttt{picxfe}) linked to slack, trend inflation expectations and relative-price shocks; foreign trade; a fiscal sector with disaggregated federal and state--local taxes, transfers and purchases, plus debt- and surplus-targeting fiscal reaction functions; and a financial block mapping the federal funds rate into term and risk-adjusted rates, equity values and other asset prices.

\subsection{Polynomial adjustment costs}

The adjustment dynamics of most key non-financial variables are governed by the generalised adjustment-cost framework known as PAC --- polynomial adjustment costs \citep{tinsley1993, tinsley2002}. An agent choosing $y_t$ minimises
\begin{equation}
C_t \;=\; \sum_{i=0}^{\infty} \beta^{i}\Big[\,\big(y_{t+i}-y^{*}_{t+i}\big)^{2} \;+\; \sum_{k=1}^{m} b_k\big((1-L)^{k}\,y_{t+i}\big)^{2}\Big],
\label{eq:pac-cost}
\end{equation}
which penalises deviations of $y$ from its desired level $y^{*}$ and movements in the level and $m-1$ higher time-differences of $y$; $\beta\;(=0.98)$ is a discount factor and $L$ the lag operator. The Euler equation can be manipulated into the PAC decision rule
\begin{equation}
\Delta y_t \;=\; a_0\big(y^{1*}_{t-1}-y_{t-1}\big) \;+\; \sum_{i=1}^{m-1} a_i\,\Delta y_{t-i} \;+\; \sum_{i=0}^{\infty} d_i\,\Delta y^{1*}_{t+i} \;+\; \sum_{i=0}^{\infty} h_i\,y^{0*}_{t+i},
\label{eq:pac-rule}
\end{equation}
where the target has been decomposed into non-stationary ($y^{1*}$) and stationary ($y^{0*}$) components: the change in $y$ responds to the lagged gap between $y$ and the trend component of its target, to its own lags (reflecting higher-order frictions), and to discounted expected future movements in the target. In the model file the two expected-future sums are assigned to explicit expectations variables $Z^{1}$ and $Z^{0}$,
\begin{equation}
\Delta y_t \;=\; a_0\big(y^{1*}_{t-1}-y_{t-1}\big) \;+\; \sum_{i=1}^{m-1} a_i\,\Delta y_{t-i} \;+\; Z^{1}_t \;+\; Z^{0}_t,
\label{eq:pac-z}
\end{equation}
augmented by a growth-neutrality correction so that $y=y^{1*}$ in a balanced-growth equilibrium. Most spending equations further combine optimising agents with rule-of-thumb agents whose spending tracks income or output $x$:
\begin{equation}
\Delta y_t \;=\; \gamma\Big(a_0\big(y^{1*}_{t-1}-y_{t-1}\big)+\sum_{i=1}^{m-1}a_i\,\Delta y_{t-i}+Z^{1}_t+Z^{0}_t\Big) \;+\; (1-\gamma)\,\Delta x_t .
\label{eq:pac-mix}
\end{equation}
PAC is used in the equations for consumption (\texttt{eco}, \texttt{ecd}), fixed investment (\texttt{eh}, \texttt{epd}, \texttt{epi}, \texttt{eps}), aggregate hours (\texttt{lhp}) and dividends, and the same error-correction-with-expectations structure organises the price--wage block.

\subsection{Expectations: VAR and MCE}

FRB/US supports two treatments of the expectations variables $Z^{0}$ and $Z^{1}$ \citep{brayton2014note}. Under \emph{model-consistent expectations} (MCE), the infinite forward sums are collapsed into finite-lead recursions in future $Z$'s and solved jointly with the model's own future paths. Under \emph{VAR expectations} --- the configuration implemented here --- expectations are generated by small estimated vector autoregressions. Each PAC equation's small model stacks the target components with a core three-variable VAR (output gap, inflation, funds rate) in companion form $z_t = H z_{t-1}$, and the expectations terms become linear functions of lagged observables:
\begin{equation}
Z^{0}_t = h_0' z_{t-1}, \qquad Z^{1}_t = h_1' z_{t-1},
\label{eq:var-z}
\end{equation}
where the coefficient vectors $h_0$, $h_1$ are known nonlinear functions of the PAC coefficients, the discount factor and the VAR companion matrix $H$. Because these expectations depend only on current and lagged variables, the VAR version of the model is fully backward-looking: the 284-equation system is simultaneous \emph{within} each quarter but recursive \emph{across} quarters, which is what permits the period-by-period solution strategy of Section~\ref{sec:implementation}.

\subsection{Policy rules}

Fiscal policy is closed by reaction functions that adjust taxes to stabilise either the debt-to-GDP ratio (\texttt{dfpdbt}) or the fiscal surplus ratio (\texttt{dfpsrp}); all experiments in this paper use the Fed's standard demo configuration (surplus-ratio targeting). Monetary policy can follow any of several published interest-rate rules selected by switch variables. The default in the Fed's demonstration programs, and in this paper, is the \emph{inertial Taylor rule} \texttt{rffintay}:
\begin{equation}
r^{ff}_t \;=\; \rho\, r^{ff}_{t-1} \;+\; (1-\rho)\Big[\, r^{*}_t + \bar\pi_t + 0.5\,\big(\bar\pi_t - \pi^{targ}_t\big) + 1.0\, \tilde x_t \Big] \;+\; \varepsilon_t,
\qquad \rho = 0.85,
\label{eq:taylor}
\end{equation}
where $\bar\pi_t$ is the four-quarter moving average of core PCE inflation, $\pi^{targ}$ the inflation target, $r^{*}$ the policymakers' perceived equilibrium real funds rate, $\tilde x_t$ the output gap (\texttt{xgap2}), and $\varepsilon_t$ the additive shock term (\texttt{rffintay\_aerr}) used to implement monetary-policy shocks. The coefficients (0.5 on the inflation gap, 1.0 on the output gap, smoothing 0.85) are as published in \texttt{model.xml}.
